\documentclass[11pt]{article}
\usepackage[a4paper,margin=1in]{geometry}
\usepackage{amsmath,amssymb,mathtools}
\usepackage{booktabs}
\usepackage{array}
\usepackage{enumitem}
\usepackage{hyperref}
\hypersetup{colorlinks=true,linkcolor=blue,urlcolor=blue,citecolor=blue}

\title{\texttt{volfi}\\[0.35em]\large Fast Black--Scholes implied-variance inversion\\[0.2em]\normalsize Technical documentation for v0.1.7 / optimized wing build\\[0.2em]\small Research prototype; not a mature or final release}
\author{Wolfgang Schadner\\\small University of Liechtenstein\\\small \href{mailto:wolfgang.schadner@uni.li}{wolfgang.schadner@uni.li}}
\date{May 2026\\[0.3em]\small Public technical disclosure draft}

\begin{document}
	\maketitle
	
	\begin{abstract}
		This document describes \texttt{volfi}, a C++ research prototype for fast inversion of the Black--Scholes normalized call price to implied total variance. The method is based on the inverse-Gaussian / generalized-inverse-Gaussian quantile representation in \cite{schadner2026explicit}. For nonzero log-moneyness, normalized calls are reduced exactly to the positive-moneyness out-of-the-money call side. The numerical kernel then approximates the total-variance GIG quantile by branchwise rational seeds and applies one analytic correction to the GIG-CDF residual in total-variance space. Version v0.1.7 adds a log-price wing seed for the true raw OTM-call domain. The public repository reports machine-precision volatility accuracy on the tested fixed and random grids, with median scalar timings of 51.80 ns/eval on the fixed benchmark grid and 55.48 ns/eval on the random benchmark grid. In that benchmark setup, \texttt{volfi} is about 3x to 3.5x faster than the benchmarked LetsBeRational interface. The exposed \texttt{otm\_context} is an implementation-level specialization for repeated inversions at fixed moneyness; analogous caching could in principle be added to other implementations.
	\end{abstract}
	
	\section{Purpose and Scope}
	
	The name of the method and implementation is \texttt{volfi}. Version identifiers identify the state of the research prototype at the time of disclosure; they are not intended to indicate a final or stable scientific version.
	
	The purpose of \texttt{volfi} is fast and accurate inversion of the Black--Scholes normalized call price. The implementation targets implied total variance rather than implied volatility itself. Let
	\begin{equation}
		k = \log\frac{K}{F}, \qquad w = \sigma^2 T, \qquad v=\sqrt{w},
	\end{equation}
	where $F$ is the forward, $K$ is the strike, $T$ is time to expiry, $\sigma$ is Black--Scholes volatility, $w$ is total variance, and $v$ is total volatility.
	
	The kernel documented here is a domain-specialized research kernel for normalized OTM-call inversion and repeated surface construction. It is not presented as a global replacement for all implied-volatility routines.
	
	\section{Normalized OTM Reduction}
	
	For nonzero log-moneyness, \texttt{volfi} maps the normalized call problem to the positive-moneyness OTM-call side. Define
	\begin{equation}
		h=|k|>0.
	\end{equation}
	If $k>0$, the normalized call is already OTM and
	\begin{equation}
		c_*=c, \qquad h=k.
	\end{equation}
	If $k<0$, the normalized call is ITM and is projected exactly to the corresponding OTM call by
	\begin{equation}
		c_* = 1 + e^{-k}(c-1), \qquad h=-k.
	\end{equation}
	At the forward strike, $k=0$, inversion is explicit:
	\begin{equation}
		v = 2\Phi^{-1}\left(\frac{1+c}{2}\right), \qquad w=v^2.
	\end{equation}
	Thus the nontrivial numerical kernel is needed only for $h>0$.
	
	\section{Variance-Space Quantile Representation}
	
	For $h>0$, the theoretical representation underlying the method is the variance-space quantile identity
	\begin{equation}
		w = F_{\mathrm{GIG}}^{-1}\left(c_*;\frac12,\frac14,h^2\right),
	\end{equation}
	where $F_{\mathrm{GIG}}^{-1}$ is the inverse distribution function of a generalized inverse Gaussian law under the stated parameterization. Equivalently, with
	\begin{equation}
		F_h(w)=F_{\mathrm{GIG}}\left(w;\frac12,\frac14,h^2\right),
	\end{equation}
	the inversion target is
	\begin{equation}
		F_h(w)=c_*, \qquad w=F_h^{-1}(c_*).
	\end{equation}
	
	This changes the approximation target from normalized volatility to total variance. \texttt{volfi} fits branchwise rational seeds to
	\begin{equation}
		c_* \mapsto F_h^{-1}(c_*),
	\end{equation}
	rather than to
	\begin{equation}
		c_* \mapsto \sqrt{F_h^{-1}(c_*)}.
	\end{equation}
	This is potentially useful because the square-root map is strongly curved near small variance, whereas the GIG representation is already expressed in the native variable $w$.
	
	The GIG quantile is not elementary. The implementation therefore uses branchwise rational seeds and one analytic correction to the GIG-CDF residual.
	
	\section{Precomputed OTM Context}
	
	For repeated inversion at fixed moneyness, \texttt{volfi} exposes an explicit OTM context storing moneyness-only quantities,
	\begin{equation}
		h, \qquad h^2, \qquad e^{h/2}, \qquad e^h.
	\end{equation}
	In v0.1.7, the true-OTM wing branch additionally uses moneyness-dependent coefficients for the logarithmic price transform.
	
	This separation is useful for surface construction, where many prices are inverted repeatedly on fixed strike or moneyness grids. The point is implementation-level specialization: the benchmarked \texttt{volfi} API exposes this reusable state directly. A different implementation of a volatility-space method could in principle introduce its own fixed-moneyness cache, but that is not part of the LetsBeRational interface used in the repository comparison.
	
	The minimal API shape is
	\begin{verbatim}
		#include <volfi/volfi.hpp>
		
		volfi::otm_context ctx(h);
		double w = volfi::implied_variance_otm(ctx, c_otm);
		
		double w2 = volfi::implied_variance_call_normalised(k, c);
	\end{verbatim}
	
	\section{Branchwise Seed Structure}
	
	The fast path has the form
	\begin{equation}
		(h,c_*) \longmapsto w_0 \longmapsto w_1,
	\end{equation}
	where $w_0$ is a branchwise rational approximation to the GIG variance quantile and $w_1$ is the corrected total variance.
	
	The optimized v0.1.7 build uses the following seed regimes:
	\begin{center}
		\begin{tabular}{ll}
			\toprule
			Region & Seed transform \\
			\midrule
			High-price region & rational seed directly in $c_*$ \\
			Middle region & rational seed in $\log(c_*/(1-c_*))$ \\
			Transition region & local rational seed \\
			True OTM wing & rational seed in $\log(c_*)$ \\
			\bottomrule
		\end{tabular}
	\end{center}
	
	The true-OTM wing branch is included because raw OTM calls can reach large positive moneyness in low-delta regions. For a call delta $\Delta$,
	\begin{equation}
		k = v\left(\frac{v}{2}-\Phi^{-1}(\Delta)\right).
	\end{equation}
	For $\Delta<1/2$,
	\begin{equation}
		k = v\left(\frac{v}{2}+|\Phi^{-1}(\Delta)|\right)>0,
	\end{equation}
	which can become large when $v$ is large. For example, at $\Delta=0.05$ and $v=2$,
	\begin{equation}
		k\approx 5.2897.
	\end{equation}
	
	\section{Variance-Space GIG-CDF Correction}
	
	The correction is formulated for the GIG-CDF residual
	\begin{equation}
		R_h(w)=F_h(w)-c_*.
	\end{equation}
	For the special GIG family used here, the CDF has the closed-form representation
	\begin{equation}
		F_h(w)=
		\Phi\left(-\frac{h}{\sqrt{w}}+\frac{\sqrt{w}}{2}\right)
		-e^h\Phi\left(-\frac{h}{\sqrt{w}}-\frac{\sqrt{w}}{2}\right).
	\end{equation}
	This representation is used as a fast evaluator of the GIG-CDF residual and its derivatives. The conceptual target remains the GIG quantile equation $F_h(w)=c_*$.
	
	Let
	\begin{equation}
		f_h(w)=F_h'(w), \qquad \ell_h(w)=\frac{f_h'(w)}{f_h(w)}.
	\end{equation}
	For this parameter family,
	\begin{equation}
		f_h(w)=\frac{e^{h/2}}{2\sqrt{2\pi w}}\exp\left(-\frac{w}{8}-\frac{h^2}{2w}\right),
	\end{equation}
	and hence
	\begin{equation}
		\ell_h(w)= -\frac{1}{2w}-\frac18+\frac{h^2}{2w^2}.
	\end{equation}
	The one-step Halley correction can be written compactly as
	\begin{equation}
		w_1=w_0-\frac{2R_h(w_0)}{2f_h(w_0)-R_h(w_0)\ell_h(w_0)}.
	\end{equation}
	The accepted optimized v0.1.7 wing build uses one such correction. No second correction step is used in the accepted wing build.
	
	\section{Relation to LetsBeRational}
	
	LetsBeRational is a highly optimized implied-volatility method by Jaeckel \cite{jaeckel2015lets}. The public \texttt{volfi} comparison uses LetsBeRational through its native \texttt{NormalisedImpliedBlackVolatility} interface. In the benchmark, the OTM call price is converted to
	\begin{equation}
		\beta=\frac{c_*}{e^{h/2}},
	\end{equation}
	and LetsBeRational is called with log-moneyness argument $-h$ and unit expiry scale. The returned total volatility is then squared to obtain total variance.
	
	The distinction is the inversion object. LetsBeRational solves the normalized implied-volatility equation in the volatility variable. \texttt{volfi} rewrites the OTM problem as a GIG quantile problem in total variance and approximates that quantile directly. Its fitted target is $F_h^{-1}(c_*)$, not $\sqrt{F_h^{-1}(c_*)}$, and its correction step is applied to $F_h(w)-c_*$ in total-variance space.
	
	\begin{center}
		\begin{tabular}{p{0.27\linewidth}p{0.31\linewidth}p{0.31\linewidth}}
			\toprule
			Aspect & \texttt{volfi} & LetsBeRational benchmark interface \\
			\midrule
			Primary unknown & total variance $w$ & total volatility $s$ \\
			Organizing identity & GIG quantile $F_h^{-1}(c_*)$ & normalized implied-volatility inversion \\
			Correction target & GIG-CDF residual $F_h(w)-c_*$ & normalized price residual in volatility space \\
			Domain specialization & normalized OTM-call side after exact projection & general normalized Black implied-volatility interface \\
			Reusable state & explicit \texttt{otm\_context} for fixed $h$ & no analogous context in the benchmark call \\
			Output in benchmark & total variance directly & total volatility, then squared \\
			\bottomrule
		\end{tabular}
	\end{center}
	
	The reusable-context row is an API and benchmark-harness distinction, not a mathematical impossibility result. Comparable fixed-moneyness caching could in principle be added to another implementation. In the public comparison, however, \texttt{volfi} is benchmarked through its precomputed \texttt{otm\_context} path, whereas LetsBeRational is called through its standard normalized implied-volatility interface.
	
	Thus, \texttt{volfi} should be described as a domain-specialized implied-variance kernel based on direct approximation and correction of a GIG quantile. Its speed advantage in the reported benchmark is obtained on the tested normalized OTM grids and under the specific comparison harness; it should not be interpreted as a universal dominance claim over LetsBeRational on every possible input domain or calling convention.
	
	\section{Validation and Benchmark Domains}
	
	This section records the domains represented in the public v0.1.7 repository. The standard \texttt{make test} target runs the fixed OTM grid, true-OTM grid, edge grid, random-delta tests, random volatility/delta tests, random true-OTM tests, and ATM test. The standalone \texttt{make bench} target runs the fixed OTM benchmark. The LetsBeRational comparison is a separate manual target because it requires a local LetsBeRational build.
	
	The fixed OTM grid used by \texttt{test\_otm\_grid}, \texttt{bench\_otm\_grid}, and the fixed LetsBeRational comparison is
	\begin{equation}
		\begin{aligned}
			v &\in \{0.01,0.05,0.10,\ldots,2.00\},\\
			\Delta &\in \{0.01,0.05,0.10,0.20,0.30,0.40,0.50,0.60,0.70,0.80,0.90\}.
		\end{aligned}
	\end{equation}
	This gives $41\times 11=451$ cases.
	
	The true-OTM fixed grid in \texttt{test\_true\_otm\_grid} uses
	\begin{equation}
		v \in \{0.01,0.05,0.10,\ldots,2.00\}, \qquad
		\Delta \in \{0.05,0.20,0.30,0.45\},
	\end{equation}
	for $41\times4=164$ cases.
	
	The random true-OTM test uses
	\begin{equation}
		v\sim U(0.01,2.0), \qquad \Delta\sim U(0.01,0.9),
	\end{equation}
	with 200,000 cases and random seed 20260502. The random projected high-delta test uses $\Delta\sim U(0.5,0.99)$ and $v\sim U(0.01,2.0)$ with the same case count and seed. The ATM test covers the same 41 volatility levels and uses the explicit forward-strike formula.
	
	\section{Accuracy Results}
	
	The public repository reports the headline accuracy comparison against LetsBeRational on the widened fixed and random grids. On the fixed 451-case grid:
	\begin{center}
		\begin{tabular}{lrrr}
			\toprule
			Method & Mean abs. vol. error & Max abs. vol. error & Max rel. vol. error \\
			\midrule
			\texttt{volfi} & $1.86\cdot 10^{-16}$ & $1.33\cdot 10^{-15}$ & $5.46\cdot 10^{-14}$ \\
			LetsBeRational & $1.64\cdot 10^{-16}$ & $7.77\cdot 10^{-16}$ & $2.32\cdot 10^{-14}$ \\
			\bottomrule
		\end{tabular}
	\end{center}
	
	On the random comparison grid:
	\begin{center}
		\begin{tabular}{lrrr}
			\toprule
			Method & Mean abs. vol. error & Max abs. vol. error & Max rel. vol. error \\
			\midrule
			\texttt{volfi} & $1.60\cdot 10^{-16}$ & $1.55\cdot 10^{-15}$ & $5.01\cdot 10^{-14}$ \\
			LetsBeRational & $1.72\cdot 10^{-16}$ & $1.33\cdot 10^{-15}$ & $3.68\cdot 10^{-14}$ \\
			\bottomrule
		\end{tabular}
	\end{center}
	
	These figures should be read as benchmark-domain results, not as a claim of global superiority on every possible normalized Black input.
	
	\section{Speed Results}
	
	The public repository reports timings for the widened fixed and random benchmark grids. These are the headline timing figures used in this document.
	
	The fixed-grid benchmark uses 451 cases, 5000 repetitions per timing run, 2,255,000 evaluations per timing run, and 9 timing runs. The reported timings are:
	\begin{center}
		\begin{tabular}{lrrrr}
			\toprule
			Method & Mean ns/eval & Median ns/eval & Min ns/eval & Max ns/eval \\
			\midrule
			\texttt{volfi} & 53.01 & 51.80 & 49.08 & 60.41 \\
			LetsBeRational & 188.10 & 181.23 & 176.99 & 211.33 \\
			\bottomrule
		\end{tabular}
	\end{center}
	The median speed ratio is
	\begin{equation}
		\frac{181.23}{51.80}\approx 3.50.
	\end{equation}
	
	The random-grid benchmark uses 200,000 accuracy cases, random seed 20260502, 5000 timing cases, 1000 repetitions per timing run, 5,000,000 evaluations per timing run, and 9 timing runs. The reported timings are:
	\begin{center}
		\begin{tabular}{lrrrr}
			\toprule
			Method & Mean ns/eval & Median ns/eval & Min ns/eval & Max ns/eval \\
			\midrule
			\texttt{volfi} & 55.88 & 55.48 & 54.19 & 59.38 \\
			LetsBeRational & 168.77 & 166.35 & 164.75 & 179.06 \\
			\bottomrule
		\end{tabular}
	\end{center}
	The median speed ratio is
	\begin{equation}
		\frac{166.35}{55.48}\approx 3.00.
	\end{equation}
	
	The benchmarked \texttt{volfi} path uses precomputed \texttt{otm\_context} values. The LetsBeRational comparison uses its standard \texttt{NormalisedImpliedBlackVolatility} interface, transforms the normalized OTM call price to the corresponding normalized input, receives total volatility, and squares it for comparison with total variance.
	
	\section{Current Bottleneck}
	
	The remaining optimization target is scalar latency in the wing region. The final residual evaluation still uses two normal CDF evaluations, one exponential-density term, one square root, and inverse-variance algebra. Consequently, further optimization should not only refit the seed. Relevant directions include:
	\begin{enumerate}[label=(\arabic*)]
		\item a stronger deep-wing seed permitting Newton-only or seed-only evaluation in subregions;
		\item a wing-specific residual approximation using Mills-ratio tail structure;
		\item finer splitting of the true-OTM wing into asymptotic regimes;
		\item a batch or SIMD API that improves throughput rather than scalar latency.
	\end{enumerate}
	
	\section{Next Research Direction}
	
	For the deep OTM wing, a plausible direction is to split the region by price and tail severity. Useful variables include
	\begin{equation}
		x=-\log(c_*), \qquad h, \qquad \frac{x}{h}, \qquad \frac{h^2}{x}, \qquad u=\frac{h}{\sqrt{w}}.
	\end{equation}
	A possible design is:
	\begin{center}
		\begin{tabular}{p{0.22\linewidth}p{0.37\linewidth}p{0.31\linewidth}}
			\toprule
			Region & Method & Purpose \\
			\midrule
			Moderate OTM & current seed + GIG-CDF correction & preserve robustness \\
			Deep OTM & asymptotic wing seed + Newton or seed-only & reduce residual cost \\
			Transition & current wing seed + GIG-CDF correction & maintain safety across boundaries \\
			\bottomrule
		\end{tabular}
	\end{center}
	
	If correction remains necessary in the deep wing, a major speed lever is to replace the two normal CDF calls inside the GIG-CDF evaluator by a restricted Mills-ratio approximation. For large $x>0$,
	\begin{equation}
		\Phi(-x)\approx \varphi(x)R(x),
	\end{equation}
	where $R(x)$ is a rational approximation to the Mills ratio.
	
	\section{Implementation Status}
	
	The public repository describes \texttt{volfi} v0.1.7 as a source-available C++ research prototype for fast Black--Scholes implied variance on the normalized OTM-call side. The repository layout includes the flagship header \texttt{include/volfi/volfi.hpp}, alternative ordering helpers, log-$c$ wing support, tests, benchmarks, and an optional LetsBeRational comparison driver. The repository does not include the LetsBeRational source or binary; that dependency must be obtained and linked separately for side-by-side comparison.
	
	The project status should be understood as a research prototype. Timings depend on machine, compiler, libm implementation, host scheduling, and benchmark harness. The reported comparison is domain-specific.
	
	\section{Technical Contribution}
	
	The documented contribution is the combination of:
	\begin{enumerate}[label=(\roman*)]
		\item exact reduction of normalized calls to the positive-moneyness OTM-call inversion problem;
		\item use of the variance-space quantile representation of \cite{schadner2026explicit} as the conceptual inversion target;
		\item branchwise rational seeds specialized to price and moneyness regimes;
		\item API-level precomputed moneyness context for repeated inversion;
		\item one analytic correction to the variance-space GIG-CDF residual;
		\item a v0.1.7 true-OTM wing branch based on a log-price transform.
	\end{enumerate}
	
	This note is intended to establish a clear technical record of the algorithmic design, tested domains, accuracy behavior, scalar timing behavior, and remaining optimization targets.
	
	\begin{thebibliography}{9}
		
		\bibitem{schadner2026explicit}
		Wolfgang Schadner.
		\newblock \emph{An Explicit Solution to Black--Scholes Implied Volatility}.
		\newblock SSRN working paper, 2026.
		\newblock Available at \url{https://papers.ssrn.com/sol3/papers.cfm?abstract_id=6649499}.
		
		\bibitem{jaeckel2015lets}
		Peter Jaeckel.
		\newblock \emph{LetsBeRational}.
		\newblock Source distribution, 2015.
		\newblock Available at \url{http://www.jaeckel.org/LetsBeRational.7z}.
		
		\bibitem{volfiRepo}
		Wolfgang Schadner.
		\newblock \emph{volfi} public repository.
		\newblock Available at \url{https://github.com/wol-fi/volfi/tree/main}.
		
	\end{thebibliography}
	
\end{document}
